% ============================================================
% chapters/ch03_allocation.tex —— 第 3 章 动态滚动的多机巢多无人机协同任务分配方法
% 小节框架已按既有方法论（task_allocation/Chapter3_README.md）拆好，
% 请逐节补充正文、公式与图示。
% ============================================================
\chapter{动态滚动的多机巢多无人机协同任务分配方法}

\section{引言}
% 本章面向广域输电网动态运维巡检，研究多机巢起降条件下的多无人机动态协同任务规划问题。

\section{问题特征与核心挑战}
\subsection{任务动态到达与任务持续滞留}
\subsection{无人机可用状态动态变化}
\subsection{多机巢起降与终点机巢选择}
\subsection{任务数量与当前执行能力不匹配}
\subsection{任务价值与空间执行代价并存}

\section{总体求解框架}
\subsection{第一阶段：动态任务选择}
\subsection{第二阶段：选择性 MUAS}

\section{第一阶段：动态任务选择机制}
\subsection{任务选择的基本思想}
\subsection{基础任务优先级}
\subsection{候选任务预筛选}
\subsection{边际收益增量选择}
\subsection{快速可行性检查}

\section{第二阶段：选择性多机巢 MUAS 优化}
\subsection{问题定义}
\subsection{决策变量}
\subsection{多目标优化模型}

\section{改进离散映射差分进化算法}
\subsection{面向任务序列的统一编码}
\subsection{多机巢终点选择}
\subsection{选择性任务分配}
\subsection{约束处理}
\subsection{三维地理代价嵌入}

\section{动态滚动与闭环运行}

\section{本章研究内容与贡献定位}
\subsection{构建动态任务选择机制}
\subsection{构建选择性多机巢 MUAS 模型及改进求解方法}
\subsection{形成动态任务选择与精细 MUAS 相结合的分层求解框架}

\section{本章总体技术路线}

\section{本章小结}

% 提示：算法示例
% \begin{algorithm}[htbp]
%   \caption{选择性多机巢 MUAS 求解}\label{alg:muas}
%   \begin{algorithmic}[1]
%     \Require 候选任务集 $\mathcal{T}_t^{\mathrm{sel}}$，可用无人机集 $U_t^{\mathrm{avail}}$
%     \Ensure 分配方案 $\mathcal{T}_t^{\mathrm{exec}}, U_t^{\mathrm{exec}}$
%     \For{每个滚动周期 $t$}
%       \State 动态任务选择 $\rightarrow \mathcal{T}_t^{\mathrm{sel}}$
%       \State 改进离散映射差分进化求解 MUAS
%     \EndFor
%     \Return 最优分配与航次
%   \end{algorithmic}
% \end{algorithm}
